\documentclass[11pt, a4paper]{article}

% --- UNIVERSAL PREAMBLE BLOCK ---
\usepackage[a4paper, top=2.5cm, bottom=2.5cm, left=2cm, right=2cm]{geometry}
\usepackage{fontspec}

\usepackage[english, bidi=basic, provide=*]{babel}

\babelprovide[import, onchar=ids fonts]{english}


% Set default/Latin font to Sans Serif in the main (rm) slot
\babelfont{rm}{Noto Sans}

% --- ADDITIONAL PACKAGES ---
\usepackage{amssymb}    % For checkboxes \square
\usepackage{booktabs}   % For professional tables
\usepackage{longtable}  % For multi-page tracking table
\usepackage{enumitem}   % For customized lists
\usepackage{xcolor}     % For colored highlights
\usepackage[colorlinks=true, linkcolor=blue]{hyperref} % For internal links

% --- CUSTOM COMMANDS ---
\newcommand{\project}[1]{\vspace{1mm}\noindent\textbf{\textcolor{blue}{Hands-on Project:}} #1}
\newcommand{\usecase}[1]{\vspace{1mm}\noindent\textbf{\textcolor{teal}{Use Case \& Architecture:}} #1}
\newcommand{\capstone}[2]{
    \vspace{3mm}
    \noindent\fbox{
        \parbox{0.97\textwidth}{
            \textbf{\textcolor{red}{CAPSTONE PROJECT - #1:}} #2
        }
    }
    \vspace{3mm}
}
\newcommand{\dayitem}[2]{\item \textbf{#1:} #2}

\title{\Huge \textbf{GCP Master Data Engineer}\\ \Large 6-Month Exhaustive Training \& Certification Program}
\author{Advanced Architecture, Hands-on Labs, and MLOps}
\date{}

\begin{document}

\maketitle

\begin{abstract}
\noindent This 24-week program is designed to transform you from a practitioner to a Master Data Engineer on Google Cloud. It progressively builds from core storage to advanced Data Mesh governance, integrating daily theory, weekly service-specific mini-projects, architectural use case analyses, and massive monthly Capstone projects. A tracking timetable is included at the end.
\end{abstract}

\tableofcontents
\newpage

% ==========================================
% MILESTONE 1
% ==========================================
\section{Milestone 1: Core Data Infrastructure \& Storage (Month 1)}
\textit{Objective: Master the foundational storage options and transactional databases.}

\subsection*{Week 1: Cloud Storage \& Transfer Services}
\begin{itemize}[label=-]
    \dayitem{Monday}{Object storage fundamentals, Storage Classes (Standard, Nearline, Coldline, Archive).}
    \dayitem{Tuesday}{Lifecycle Management, Versioning, and Retention Policies.}
    \dayitem{Wednesday}{Data Migration: Storage Transfer Service vs. Transfer Appliance.}
    \dayitem{Thursday}{\project{Deploy a GCS bucket via CLI, configure a 30-day transition to Archive, and set up a Storage Transfer job from a public AWS S3 bucket.}}
    \dayitem{Friday}{\usecase{Design a multi-region disaster recovery storage architecture for a media streaming company serving 10PB of video assets.}}
\end{itemize}

\subsection*{Week 2: Relational Databases (Cloud SQL \& AlloyDB)}
\begin{itemize}[label=-]
    \dayitem{Monday}{Cloud SQL Architecture, High Availability (HA), and Read Replicas.}
    \dayitem{Tuesday}{Database Migration Service (DMS) and Private IP connections.}
    \dayitem{Wednesday}{AlloyDB for PostgreSQL: Columnar engine and ML integration.}
    \dayitem{Thursday}{\project{Deploy an HA Cloud SQL Postgres instance, create a read replica, and simulate a regional failover.}}
    \dayitem{Friday}{\usecase{Architect a zero-downtime migration strategy for a monolithic on-premise ERP system to Cloud SQL.}}
\end{itemize}

\subsection*{Week 3: Global Relational Scaling (Cloud Spanner)}
\begin{itemize}[label=-]
    \dayitem{Monday}{Spanner concepts: Nodes, Replicas, TrueTime, and Strong Global Consistency.}
    \dayitem{Tuesday}{Schema Design: Primary keys, Interleaved Tables, and Indexing.}
    \dayitem{Wednesday}{Spanner performance tuning and query execution plans.}
    \dayitem{Thursday}{\project{Create a multi-region Spanner instance, design an interleaved schema for an e-commerce cart, and load 10,000 synthetic records.}}
    \dayitem{Friday}{\usecase{Design a globally distributed payment processing ledger that requires 99.999\% availability and zero data loss across continents.}}
\end{itemize}

\subsection*{Week 4: Wide-Column NoSQL (Cloud Bigtable)}
\begin{itemize}[label=-]
    \dayitem{Monday}{Bigtable architecture: Tablets, Nodes, and App Profiles.}
    \dayitem{Tuesday}{Row Key Design: Salting, hashing, and avoiding hotspotting.}
    \dayitem{Wednesday}{Performance optimization, garbage collection, and Key Visualizer.}
    \dayitem{Thursday}{\project{Deploy a Bigtable cluster, write a Python script to ingest IoT temperature sensor data using a reversed-timestamp row key.}}
    \dayitem{Friday}{\usecase{Design an ingestion and serving layer for a financial trading platform processing 100,000 stock ticks per second.}}
\end{itemize}

\capstone{Milestone 1}{Build an E-Commerce Multi-Tier Data Backend. Deploy Cloud SQL for user accounts, Spanner for global inventory/checkout transactions, Bigtable for clickstream tracking, and Cloud Storage for product images. Map the VPC peering required to connect them securely.}

% ==========================================
% MILESTONE 2
% ==========================================
\section{Milestone 2: Data Warehousing \& Analytics (Month 2)}
\textit{Objective: Deep dive into BigQuery optimization, Lakehouses, and Data Sharing.}

\subsection*{Week 5: BigQuery Architecture \& Core SQL}
\begin{itemize}[label=-]
    \dayitem{Monday}{BigQuery decoupled architecture: Colossus, Jupiter, Dremel.}
    \dayitem{Tuesday}{Pricing models (On-demand vs. Capacity/Editions) and Slots.}
    \dayitem{Wednesday}{Advanced SQL: Window functions, Arrays, and Structs.}
    \dayitem{Thursday}{\project{Query the GitHub public dataset, unnest arrays, and calculate a 7-day rolling average of commits per repository using window functions.}}
    \dayitem{Friday}{\usecase{Estimate the monthly cost of a marketing analytics workload scanning 50TB daily and recommend an Edition tier.}}
\end{itemize}

\subsection*{Week 6: BigQuery Optimization}
\begin{itemize}[label=-]
    \dayitem{Monday}{Partitioning: Ingestion-time vs. Column-based (Time/Integer).}
    \dayitem{Tuesday}{Clustering: Sorting data for optimal filter performance.}
    \dayitem{Wednesday}{Materialized Views, BI Engine, and Query execution plan analysis.}
    \dayitem{Thursday}{\project{Create a massive synthetic table. Compare the execution time and bytes billed between a full scan, a partitioned scan, and a clustered scan.}}
    \dayitem{Friday}{\usecase{Re-architect a poorly performing multi-petabyte log analysis table that is currently causing high query costs for the security team.}}
\end{itemize}

\subsection*{Week 7: Advanced Data Types \& Federated Queries}
\begin{itemize}[label=-]
    \dayitem{Monday}{BigQuery GIS (Geospatial) and BigQuery Search.}
    \dayitem{Tuesday}{External Tables and querying data directly in Cloud Storage.}
    \dayitem{Wednesday}{BigQuery Omni (Cross-cloud analytics on AWS/Azure).}
    \dayitem{Thursday}{\project{Create an external table pointing to CSVs in GCS. Write a query joining this external data with native BigQuery GIS taxi-dropoff data.}}
    \dayitem{Friday}{\usecase{Design a multi-cloud data architecture for a retail company that has historical sales in AWS S3 but wants to analyze them in GCP.}}
\end{itemize}

\subsection*{Week 8: Data Lakehouse \& Analytics Hub}
\begin{itemize}[label=-]
    \dayitem{Monday}{BigLake: Unifying data lakes and data warehouses with fine-grained security.}
    \dayitem{Tuesday}{Object Tables: Analyzing unstructured data (PDFs/Images).}
    \dayitem{Wednesday}{Analytics Hub: Secure data sharing and Data Clean Rooms.}
    \dayitem{Thursday}{\project{Set up a BigLake table with column-level access control. Publish a dataset to the Analytics Hub for a "partner" organization.}}
    \dayitem{Friday}{\usecase{Architect a Data Clean Room for an ad agency to merge their CRM data with Google Ads data without exposing raw PII.}}
\end{itemize}

\capstone{Milestone 2}{Enterprise Data Lakehouse Build. Ingest 5 years of retail data into GCS. Create BigLake tables over the data, apply row/column level security policies based on region, create Materialized Views for executive dashboards, and set up a BI Engine reservation.}

% ==========================================
% MILESTONE 3
% ==========================================
\section{Milestone 3: Batch Processing \& ETL Orchestration (Month 3)}
\textit{Objective: Master Hadoop migration, No-code ETL, CDC, and Airflow.}

\subsection*{Week 9: Managed Hadoop/Spark (Cloud Dataproc)}
\begin{itemize}[label=-]
    \dayitem{Monday}{Dataproc Architecture, Ephemeral clusters, and Preemptible VMs.}
    \dayitem{Tuesday}{Migrating from on-prem HDFS to Cloud Storage.}
    \dayitem{Wednesday}{Serverless Dataproc and Spark job tuning.}
    \dayitem{Thursday}{\project{Deploy an ephemeral Dataproc cluster, submit a PySpark job that reads from GCS, aggregates sales data, writes back to GCS, and self-deletes.}}
    \dayitem{Friday}{\usecase{Design a migration strategy for a 100-node on-premise Cloudera cluster, prioritizing cost reduction and decoupling compute from storage.}}
\end{itemize}

\subsection*{Week 10: Visual ETL (Cloud Data Fusion)}
\begin{itemize}[label=-]
    \dayitem{Monday}{CDAP concepts, Data Fusion architecture, and Private IPs.}
    \dayitem{Tuesday}{Wrangler: Interactive data cleaning and transformation.}
    \dayitem{Wednesday}{Plugins, Macros, and Scheduling pipelines.}
    \dayitem{Thursday}{\project{Build a no-code Data Fusion pipeline that reads messy customer data from GCS, uses Wrangler to extract domains from emails, and writes to BigQuery.}}
    \dayitem{Friday}{\usecase{Evaluate Data Fusion vs. Dataflow for a team of business analysts who need to build daily reporting pipelines without writing Java/Python.}}
\end{itemize}

\subsection*{Week 11: Change Data Capture (Datastream)}
\begin{itemize}[label=-]
    \dayitem{Monday}{CDC concepts and the Datastream architecture.}
    \dayitem{Tuesday}{Source configurations (MySQL/PostgreSQL/Oracle).}
    \dayitem{Wednesday}{Destination configurations (Cloud Storage vs. BigQuery native).}
    \dayitem{Thursday}{\project{Set up a Cloud SQL PostgreSQL database, configure logical replication, and build a Datastream pipeline to sync changes to BigQuery in near real-time.}}
    \dayitem{Friday}{\usecase{Architect a zero-downtime database replication pipeline from an on-premise Oracle database to BigQuery for real-time BI reporting.}}
\end{itemize}

\subsection*{Week 12: Pipeline Orchestration (Cloud Composer)}
\begin{itemize}[label=-]
    \dayitem{Monday}{Apache Airflow concepts: DAGs, Operators, Tasks, and XComs.}
    \dayitem{Tuesday}{Cloud Composer 2 architecture (Autoscaling, GKE underlying).}
    \dayitem{Wednesday}{Cross-service orchestration (Triggering Dataproc, Dataflow, BQ).}
    \dayitem{Thursday}{\project{Write a Python DAG that: 1. Checks if a file exists in GCS. 2. Triggers a Dataproc job. 3. Loads the output to BigQuery. 4. Sends an email on failure.}}
    \dayitem{Friday}{\usecase{Design a resilient, self-healing orchestration layer for a daily banking reconciliation process spanning 15 dependent systems.}}
\end{itemize}

\capstone{Milestone 3}{Automated Batch ELT Pipeline. Use Cloud Scheduler to trigger a Composer DAG. The DAG runs a Datastream CDC sync verification, triggers a Serverless Dataproc Spark job to calculate complex ML features, loads them into BigQuery, and creates a Materialized View.}

\newpage
% ==========================================
% MILESTONE 4
% ==========================================
\section{Milestone 4: Streaming \& Real-Time Analytics (Month 4)}
\textit{Objective: Conquer unbounded data, message queues, and real-time visualization.}

\subsection*{Week 13: Asynchronous Messaging (Pub/Sub)}
\begin{itemize}[label=-]
    \dayitem{Monday}{Pub/Sub concepts: Topics, Subscriptions (Push/Pull), Acknowledgements.}
    \dayitem{Tuesday}{Dead-Letter Topics, Retry Policies, and Snapshot/Seek.}
    \dayitem{Wednesday}{Pub/Sub Lite vs. Pub/Sub Standard (Cost vs. Operational overhead).}
    \dayitem{Thursday}{\project{Create a Pub/Sub topic and pull subscription. Write a Python publisher to stream mock financial transactions. Configure a Dead-Letter Topic for failed messages.}}
    \dayitem{Friday}{\usecase{Architect a global telemetry ingestion point for 5 million mobile games, ensuring exactly-once processing semantics.}}
\end{itemize}

\subsection*{Week 14: Dataflow \& Apache Beam (Basics)}
\begin{itemize}[label=-]
    \dayitem{Monday}{Apache Beam model: PCollections, PTransforms, ParDo, and Runners.}
    \dayitem{Tuesday}{Dataflow architecture: Streaming Engine, Shuffle, Autoscaling.}
    \dayitem{Wednesday}{Dataflow Templates (Classic vs. Flex Templates).}
    \dayitem{Thursday}{\project{Write an Apache Beam pipeline in Python that reads texts, tokenizes words, counts frequencies (WordCount), and runs it on the Dataflow runner.}}
    \dayitem{Friday}{\usecase{Design a pipeline to anonymize incoming healthcare HL7 messages in real-time before storing them in BigQuery.}}
\end{itemize}

\subsection*{Week 15: Dataflow Advanced (Windowing \& Time)}
\begin{itemize}[label=-]
    \dayitem{Monday}{Event Time vs. Processing Time and Watermarks.}
    \dayitem{Tuesday}{Windowing strategies: Fixed, Sliding, and Session windows.}
    \dayitem{Wednesday}{Handling Late Data: Allowed Lateness and Triggers.}
    \dayitem{Thursday}{\project{Build a streaming Dataflow pipeline reading from Pub/Sub. Group user clicks into 5-minute Session Windows, allowing for 2 minutes of late data.}}
    \dayitem{Friday}{\usecase{Troubleshoot a scenario where a streaming dashboard is showing inaccurate daily aggregations due to mobile devices going offline and uploading data late.}}
\end{itemize}

\subsection*{Week 16: Business Intelligence \& Looker}
\begin{itemize}[label=-]
    \dayitem{Monday}{Looker Architecture: The Semantic Layer and LookML.}
    \dayitem{Tuesday}{Dimensions, Measures, Explores, and Views.}
    \dayitem{Wednesday}{Looker Studio vs. Looker Core (Self-serve vs. Governed BI).}
    \dayitem{Thursday}{\project{Connect Looker Studio (or a Looker trial) to your BigQuery retail dataset. Create a real-time dashboard monitoring streaming sales with drill-down filters.}}
    \dayitem{Friday}{\usecase{Design the BI access layer for an enterprise where 500 analysts need to query centralized data without creating conflicting metrics (e.g., "Gross Revenue" definitions).}}
\end{itemize}

\capstone{Milestone 4}{Real-Time Fraud Detection. Build a Python script simulating credit card swipes sent to Pub/Sub. Create a Dataflow pipeline that uses a Sliding Window to detect if a card is used in two different cities within 10 minutes, flags it, and streams the alert to BigQuery for live Looker visualization.}


% ==========================================
% MILESTONE 5
% ==========================================
\section{Milestone 5: Machine Learning Operationalization (Month 5)}
\textit{Objective: Embed AI into data pipelines and manage the MLOps lifecycle.}

\subsection*{Week 17: BigQuery ML (BQML)}
\begin{itemize}[label=-]
    \dayitem{Monday}{BQML capabilities: Regression, Classification, K-Means Clustering.}
    \dayitem{Tuesday}{Time-series forecasting (ARIMA\_PLUS) and Anomaly Detection.}
    \dayitem{Wednesday}{Model evaluation (ML.EVALUATE) and hyperparameter tuning.}
    \dayitem{Thursday}{\project{Use BQML to train a logistic regression model on Google Analytics sample data to predict user churn. Evaluate the model's ROC curve using SQL.}}
    \dayitem{Friday}{\usecase{Architect a system to automatically forecast daily inventory needs for 1,000 retail stores using purely SQL-based pipelines.}}
\end{itemize}

\subsection*{Week 18: Vertex AI Foundations}
\begin{itemize}[label=-]
    \dayitem{Monday}{Vertex AI Ecosystem, AutoML vs. Custom Training.}
    \dayitem{Tuesday}{Vertex AI Model Registry and Endpoint deployment (Online vs. Batch).}
    \dayitem{Wednesday}{Explainable AI (Feature attributions).}
    \dayitem{Thursday}{\project{Train an AutoML Tabular model on a cleaned BigQuery dataset. Deploy the model to a Vertex AI Endpoint and send a cURL request for an online prediction.}}
    \dayitem{Friday}{\usecase{Design a workflow for data scientists to train custom TensorFlow models on GPUs while accessing data securely stored in BigQuery.}}
\end{itemize}

\subsection*{Week 19: MLOps \& Vertex AI Pipelines}
\begin{itemize}[label=-]
    \dayitem{Monday}{Vertex AI Feature Store: Online vs. Offline serving.}
    \dayitem{Tuesday}{Kubeflow pipelines and Vertex AI Pipelines orchestration.}
    \dayitem{Wednesday}{Continuous Training (CT) triggers and Model Monitoring (Data Drift).}
    \dayitem{Thursday}{\project{Build a simple Kubeflow pipeline (Extract $\rightarrow$ Train $\rightarrow$ Deploy). Compile it and run it using Vertex AI Pipelines.}}
    \dayitem{Friday}{\usecase{Architect an MLOps lifecycle that detects data drift in a production recommendation engine and automatically triggers a retraining pipeline via Cloud Functions.}}
\end{itemize}

\subsection*{Week 20: Unstructured Data \& AI APIs}
\begin{itemize}[label=-]
    \dayitem{Monday}{Cloud Vision API and Video Intelligence API.}
    \dayitem{Tuesday}{Natural Language API and Translation API.}
    \dayitem{Wednesday}{Integrating AI APIs within Dataflow pipelines.}
    \dayitem{Thursday}{\project{Write a Dataflow pipeline that reads customer support text reviews from Pub/Sub, calls the Natural Language API to get sentiment scores, and saves to BigQuery.}}
    \dayitem{Friday}{\usecase{Design a scalable pipeline to process terabytes of uploaded user images, extract text (OCR), tag objects, and make the metadata searchable.}}
\end{itemize}

\capstone{Milestone 5}{End-to-End MLOps Pipeline. Use Composer to orchestrate a pipeline that extracts features from BigQuery, pushes them to Vertex Feature Store, trains a custom model using Vertex Training, registers the model, and sets up Vertex Model Monitoring for serving skew.}

\newpage
% ==========================================
% MILESTONE 6
% ==========================================
\section{Milestone 6: Security, Governance \& Exam Prep (Month 6)}
\textit{Objective: Secure the perimeter, implement Data Mesh, and certify.}

\subsection*{Week 21: Security, IAM \& Networking}
\begin{itemize}[label=-]
    \dayitem{Monday}{VPC Service Controls (Perimeters) and Private Service Connect.}
    \dayitem{Tuesday}{IAM Deep Dive: Custom Roles, Service Accounts, and Workload Identity.}
    \dayitem{Wednesday}{Cloud KMS (CMEK vs. CSEK) in BigQuery and GCS.}
    \dayitem{Thursday}{\project{Create a VPC Service Perimeter around a BigQuery dataset. Attempt to query it from an external IP to verify the block. Configure an access level to allow your IP.}}
    \dayitem{Friday}{\usecase{Design a highly secure data perimeter for a healthcare provider ensuring no data exfiltration can occur even if IAM credentials are compromised.}}
\end{itemize}

\subsection*{Week 22: Data Governance \& Data Mesh (Dataplex)}
\begin{itemize}[label=-]
    \dayitem{Monday}{Dataplex architecture: Lakes, Zones (Raw/Curated), and Assets.}
    \dayitem{Tuesday}{Data Catalog: Tag templates, Search, and Data Lineage.}
    \dayitem{Wednesday}{Auto Data Quality (AutoDQ) and Data Profiling.}
    \dayitem{Thursday}{\project{Set up a Dataplex Lake. Map a GCS bucket and BigQuery dataset as assets. Create a tag template and apply policy tags to restrict access to a 'Salary' column.}}
    \dayitem{Friday}{\usecase{Architect a Data Mesh for a global conglomerate with 5 independent domains, ensuring centralized governance but decentralized data ownership.}}
\end{itemize}

\subsection*{Week 23: Compliance \& Data Loss Prevention (DLP)}
\begin{itemize}[label=-]
    \dayitem{Monday}{Cloud DLP concepts: InfoTypes, Inspection, and De-identification.}
    \dayitem{Tuesday}{Masking, Tokenization, and Crypto-based hashing.}
    \dayitem{Wednesday}{Automated DLP scanning in BigQuery and GCS.}
    \dayitem{Thursday}{\project{Write a Python script using the DLP API to scan a block of text, identify emails and credit card numbers, and replace them with asterisks (masking).}}
    \dayitem{Friday}{\usecase{Design a pipeline that ingests third-party customer data, automatically identifies PII, tokenizes it using a KMS key, and stores the safe version in the data warehouse.}}
\end{itemize}

\subsection*{Week 24: Final Certification Preparation}
\begin{itemize}[label=-]
    \dayitem{Monday}{Case Study Review 1: EHR Healthcare (focus on compliance and migration).}
    \dayitem{Tuesday}{Case Study Review 2: Helium Sports (focus on real-time analytics).}
    \dayitem{Wednesday}{Full-length Mock Exam 1 (Identify weak domains).}
    \dayitem{Thursday}{Targeted review of weak domains based on mock exam results.}
    \dayitem{Friday}{Full-length Mock Exam 2.}
\end{itemize}

\capstone{Milestone 6}{Secure, Governed Data Mesh. Build a multi-project architecture. Project A ingests streaming PII data and uses DLP to tokenize it. Project B (Analytics) hosts BigLake tables. Project C (Governance) uses Dataplex to enforce row-level security and VPC Service Controls to lock down all API access.}

\newpage
% ==========================================
% TRACKING TIMETABLE
% ==========================================
\section{6-Month Progress Tracking Timetable}

\begin{center}
\renewcommand{\arraystretch}{1.5}
\begin{longtable}{@{} c l p{7cm} c @{}}
\toprule
\textbf{Week} & \textbf{Primary Focus} & \textbf{Key Project / Capstone} & \textbf{Status} \\ 
\midrule
\endfirsthead

\toprule
\textbf{Week} & \textbf{Primary Focus} & \textbf{Key Project / Capstone} & \textbf{Status} \\ 
\midrule
\endhead

\bottomrule
\endfoot

\bottomrule
\endlastfoot

1 & GCS \& Transfers & GCS Lifecycle \& Transfer Job & $\square$ \\
2 & Cloud SQL \& AlloyDB & HA PostgreSQL Failover Simulation & $\square$ \\
3 & Cloud Spanner & Global Interleaved Schema & $\square$ \\
4 & Cloud Bigtable & \textbf{Cap 1: E-Commerce Multi-Tier Data} & $\square$ \\
\midrule
5 & BigQuery Architecture & Advanced SQL Window Functions & $\square$ \\
6 & BQ Optimization & Partitioning \& Clustering Benchmarks & $\square$ \\
7 & BQ Advanced & GIS \& External Tables & $\square$ \\
8 & Lakehouse \& Analytics Hub & \textbf{Cap 2: Enterprise Data Lakehouse} & $\square$ \\
\midrule
9 & Cloud Dataproc & Spark Job on Ephemeral Cluster & $\square$ \\
10 & Cloud Data Fusion & No-code Wrangler ETL Pipeline & $\square$ \\
11 & Datastream (CDC) & PostgreSQL to BQ Logical Sync & $\square$ \\
12 & Cloud Composer & \textbf{Cap 3: Automated Batch ELT} & $\square$ \\
\midrule
13 & Cloud Pub/Sub & DLQ and Push/Pull Scripting & $\square$ \\
14 & Dataflow Basics & Beam WordCount on Cloud Runner & $\square$ \\
15 & Dataflow Advanced & Session Windows \& Late Data logic & $\square$ \\
16 & Looker \& BI & \textbf{Cap 4: Real-Time Fraud Telemetry} & $\square$ \\
\midrule
17 & BigQuery ML & Churn Prediction \& Evaluation & $\square$ \\
18 & Vertex AI Foundations & AutoML Deployment \& cURL Predict & $\square$ \\
19 & MLOps \& Pipelines & Kubeflow Pipeline Orchestration & $\square$ \\
20 & AI APIs & \textbf{Cap 5: End-to-End MLOps Pipeline} & $\square$ \\
\midrule
21 & IAM \& VPC SC & Service Perimeters \& Access Levels & $\square$ \\
22 & Dataplex \& Data Mesh & AutoDQ and Policy Tags & $\square$ \\
23 & Cloud DLP & Python Tokenization/Masking Script & $\square$ \\
24 & Exam Prep & \textbf{Cap 6: Secure Data Mesh Design} & $\square$ \\

\end{longtable}
\end{center}

\vspace{1cm}
\begin{center}
    \textit{"Consistency is the architecture of mastery."} \\
    Best of luck on your 6-month journey to becoming a Master Google Cloud Data Engineer!
\end{center}

\end{document}